% LLNCS — bản thảo tiếng Việt (main_vi).
% Khung phương pháp mới: CTI/NAR/QDR theo độ cụ thể của cam kết (đã bỏ EWRI/grounding).
% Lưu ý biên dịch: tiếng Việt nên dùng XeLaTeX/LuaLaTeX (fontspec) cho dấu phụ;
% nếu dùng pdfLaTeX, giữ inputenc utf8 + fontenc T1 và một font hỗ trợ tiếng Việt.
\documentclass[runningheads]{llncs}
%
\usepackage[utf8]{inputenc}
\usepackage[T1]{fontenc}
\usepackage{graphicx}
\usepackage[colorlinks=true,linkcolor=blue,citecolor=blue]{hyperref}
\usepackage{amsmath}
\usepackage{multirow}
\usepackage{booktabs}

\begin{document}
%
\title{Đo lường rủi ro ESG-washing qua độ cụ thể của cam kết: Nghiên cứu trên báo cáo thường niên ngân hàng thương mại Việt Nam}
%
\titlerunning{Đo lường ESG-washing trong báo cáo ngân hàng Việt Nam}
%
\author{Huu Huy Pham \and Viet Anh Nguyen\thanks{Tác giả liên hệ: vietanh@vnu.edu.vn}}
%
\authorrunning{H. H. Pham and V. A. Nguyen}
%
\institute{Trường Đại học Công nghệ, Đại học Quốc gia Hà Nội, Việt Nam\\
\email{22028237@vnu.edu.vn}, \email{vietanh@vnu.edu.vn}}
%
\maketitle
%
\begin{abstract}
Công bố thông tin ESG đã trở thành cấu phần trung tâm trong báo cáo của ngành ngân hàng, nhưng sự bùng nổ của ngôn ngữ ESG cũng làm gia tăng lo ngại về \emph{ESG-washing} --- xu hướng nhấn mạnh các cam kết nghe hay nhưng thiếu nội dung thực chất. Các phương pháp phát hiện hiện có phần lớn dựa vào điểm xếp hạng ESG của bên thứ ba (vốn thưa và kém tin cậy ở thị trường mới nổi) hoặc dựa vào việc đối chiếu bằng chứng \emph{nội văn bản} --- một cấu trúc đo lường mang tính vòng tròn vì bằng chứng và cam kết đều do chính ngân hàng viết trong cùng một báo cáo. Nghiên cứu này đề xuất đo ESG-washing trực tiếp qua \emph{độ cụ thể của cam kết} (commitment specificity): mỗi cam kết ESG được một mô hình ngôn ngữ lớn chấm theo thang ba mức (mơ hồ / hành động có tên / định lượng quy về chủ thể), từ đó tổng hợp thành ba chỉ số bổ sung nhau trên mỗi (ngân hàng, năm): Tỉ lệ nói suông (CTI), Tỉ lệ hành động có tên (NAR) và Tỉ lệ công bố định lượng (QDR). Khung đo được áp dụng cho 6 ngân hàng thương mại Việt Nam trong giai đoạn 2020--2024 (30 quan sát panel). Kết quả cho thấy: (i) trung bình 36{,}6\% cam kết ESG là nói suông; (ii) trụ Môi trường bị công bố thiếu một cách có ý nghĩa thống kê so với hai trụ còn lại (kiểm định Friedman, $p < 10^{-6}$), phản ánh hành vi né chủ đề khó; và (iii) khối lượng cam kết tăng mạnh theo thời gian trong khi tỉ trọng nội dung thực chất không cải thiện --- tức ``nói nhiều hơn nhưng không cụ thể hơn''. Xếp hạng rủi ro giữa các ngân hàng ổn định dưới kiểm định bootstrap (Kendall $\tau = 0{,}67$). Nghiên cứu cung cấp một khung đo ESG-washing trung thực, đầy đủ ba trụ E/S/G, không cần dữ liệu bên thứ ba, phù hợp cho các thị trường thiếu dữ liệu.

\keywords{ESG-washing \and Báo cáo thường niên \and Ngân hàng \and Xử lý ngôn ngữ tự nhiên \and Độ cụ thể của cam kết}
\end{abstract}
%
%
\section{Giới thiệu}
Công bố thông tin ESG (Môi trường -- Xã hội -- Quản trị) đã chuyển từ hoạt động tự nguyện thành một cấu phần cốt lõi trong báo cáo của ngành ngân hàng, nơi các cam kết phát triển bền vững ngày càng được cơ quan quản lý, nhà đầu tư và công chúng giám sát chặt chẽ \cite{buchetti2025literature,lagasio2024esg}. Báo cáo thường niên là kênh chính để truyền tải các sáng kiến ESG, nhưng sự gia tăng của các diễn ngôn ESG cũng khuếch đại rủi ro \emph{ESG-washing} --- thực hành nhấn mạnh các tuyên bố thuận lợi mà thiếu bằng chứng thực chất kiểm chứng được \cite{florstedt2025detecting,delmas2011drivers}.

Về mặt lý thuyết, ESG-washing là một dạng \emph{decoupling} (tách rời) giữa \emph{lời nói} (talk) và \emph{việc làm} (walk): tổ chức phát đi tín hiệu phù hợp với kỳ vọng xã hội nhưng không thực thi tương ứng \cite{bromley2012smoke,marquis2016scrutiny}. Hệ quả là làm xói mòn niềm tin của các bên liên quan, bóp méo môi trường thông tin và dẫn tới phân bổ vốn kém hiệu quả \cite{lagasio2024esg,kathan2025you}. Những rủi ro này đặc biệt nghiêm trọng trong ngành ngân hàng, như minh hoạ qua tranh cãi năm 2022 liên quan các quỹ đầu tư bền vững của Deutsche Bank AG \cite{lipenkova2023detecting}.

Các hướng phát hiện ESG-washing hiện nay có thể chia thành hai dòng. Dòng thứ nhất dựa trên \emph{chênh lệch xếp hạng}, đo độ lệch giữa công bố của doanh nghiệp và điểm ESG của bên thứ ba \cite{su16062571,ghitti2024agency,kathan2025you}. Dòng thứ hai dựa trên \emph{phân tích văn bản} và NLP để khai thác đặc trưng ngôn ngữ trong công bố ESG \cite{lagasio2024esg,lipenkova2023detecting,zhao2023detecting,bingler2022cheap}. Tuy nhiên, cả hai dòng đều gặp trở ngại lớn ở các thị trường mới nổi như Việt Nam: dữ liệu xếp hạng và giám sát độc lập của bên thứ ba gần như không có hoặc rời rạc, trong khi phần lớn công cụ NLP được thiết kế cho tiếng Anh và chỉ tập trung vào trụ Môi trường.

Một số nghiên cứu gần đây cố gắng vượt qua việc thiếu dữ liệu ngoài bằng cách đối chiếu mỗi tuyên bố ESG với \emph{bằng chứng nội văn bản} trong cùng báo cáo (truy hồi câu hỗ trợ rồi chấm bằng suy luận ngôn ngữ) \cite{florstedt2025detecting}. Chúng tôi lập luận rằng cách tiếp cận này mang tính \emph{vòng tròn}: cả tuyên bố lẫn ``bằng chứng'' đều do chính ngân hàng viết, nên việc một câu được câu khác trong cùng tài liệu hỗ trợ chỉ xác nhận \emph{báo cáo nhất quán với chính nó}, chứ không xác nhận hành động thực tế bên ngoài. Do Việt Nam không có walk-data của bên thứ ba, chúng tôi không giả vờ đo ``walk'', mà đo một đại lượng \emph{thật sự rút được từ văn bản}: cam kết \emph{cụ thể} đến đâu.

Trên cơ sở đó, nghiên cứu đề xuất một khung NLP đo ESG-washing qua \emph{độ cụ thể của cam kết} (commitment specificity). Mỗi cam kết ESG được phân loại thành ba mức --- mơ hồ, hành động có tên, hoặc định lượng quy về chủ thể --- rồi tổng hợp thành ba chỉ số bổ sung trên mỗi (ngân hàng, năm): \textbf{CTI} (Cheap Talk Index, tỉ lệ nói suông), \textbf{NAR} (Named Action Rate) và \textbf{QDR} (Quantified Disclosure Rate). Cách đo này dựa trực tiếp vào lý thuyết ``cheap talk'' \cite{farrell1996cheap}: phát ngôn không tốn chi phí và không ràng buộc thì ít giá trị thông tin; tính cụ thể và khả năng quy trách nhiệm chính là chi phí phân biệt cam kết thực với cam kết suông.

Đóng góp của nghiên cứu gồm bốn điểm. \textbf{Thứ nhất}, đây là khung đo ESG-washing đầy đủ ba trụ E/S/G đầu tiên cho ngân hàng Việt Nam, không phụ thuộc walk-data bên ngoài. \textbf{Thứ hai}, về phương pháp, chúng tôi thay cấu trúc grounding vòng tròn bằng một thang \emph{specificity} thứ tự ba mức, có nhận biết quy-về-chủ-thể (attribution-aware) và cơ chế chống ``bịa số''. \textbf{Thứ ba}, chúng tôi cung cấp bức tranh thực nghiệm panel cùng bộ kiểm định thống kê (bootstrap CI, Friedman/Wilcoxon cho công bố chọn lọc, kiểm định ổn định xếp hạng). \textbf{Thứ tư}, toàn bộ pipeline và dữ liệu translate-train được công khai để tái lập.

Phần còn lại của bài được tổ chức như sau. Mục~2 điểm lại các nghiên cứu liên quan. Mục~3 trình bày khung phương pháp và cách xây dựng các chỉ số. Mục~4 trình bày và thảo luận kết quả thực nghiệm. Mục~5 kết luận và nêu hướng phát triển.

%
%
\section{Các nghiên cứu liên quan}
\subsection{ESG-washing và lý thuyết decoupling}
Văn liệu quản trị mô tả greenwashing/ESG-washing như sự kết hợp giữa công bố tích cực và hiệu quả môi trường--xã hội kém \cite{delmas2011drivers,lyon2015means}. Marquis và cộng sự \cite{marquis2016scrutiny} cho thấy doanh nghiệp phản ứng với áp lực giám sát bằng \emph{công bố chọn lọc} (selective disclosure): khuếch đại thông tin tích cực, che giấu thông tin bất lợi. Cho và Patten \cite{cho2007role} lý giải công bố môi trường như công cụ tạo \emph{tính chính danh}. Ở tầng tổ chức, Bromley và Powell \cite{bromley2012smoke} hệ thống hoá hiện tượng decoupling giữa chính sách tuyên bố và thực thi. Những khung lý thuyết này định hướng hai câu hỏi đo lường của chúng tôi: cam kết có \emph{cụ thể} không, và ngân hàng có \emph{né trụ khó} không.

\subsection{Phát hiện ESG-washing bằng NLP}
Florstedt và cộng sự \cite{florstedt2025detecting} đề xuất khung lai kết hợp phân tích cảm xúc (VADER), đánh giá độ cụ thể bằng ClimateBERT \cite{webersinke2021climatebert} và truy hồi bằng chứng nội văn bản bằng Sentence-BERT \cite{reimers2019sentence}; tuy nhiên khung này thiết kế cho tiếng Anh và lệ thuộc bằng chứng cùng tài liệu. Lipenkova và cộng sự \cite{lipenkova2023detecting} đối chiếu công bố của doanh nghiệp Đức với nguồn tin bên ngoài, nhưng phụ thuộc nặng vào dữ liệu ngoài --- thứ khan hiếm ở thị trường mới nổi. Lagasio \cite{lagasio2024esg} xây dựng chỉ số ESGSI dựa trên TF--IDF, LDA và phân tích cảm xúc, nắm bắt độ lệch giữa ngôn ngữ tích cực và chiều sâu nội dung, nhưng thiên về đặc trưng phong cách/chủ đề và không đánh giá tính kiểm chứng của từng tuyên bố. Bingler và cộng sự \cite{bingler2022cheap} dùng ClimateBERT để chỉ ra hiện tượng ``cheap talk và cherry-picking'' trong công bố rủi ro khí hậu --- gần với động lực của nghiên cứu này, song giới hạn ở tiếng Anh và trụ khí hậu.

\subsection{Khoảng trống nghiên cứu}
Tổng hợp lại, văn liệu hiện tồn tại ba khoảng trống. \emph{Một}, phần lớn công cụ thiết kế cho tiếng Anh và trụ Môi trường, ít áp dụng cho ngôn ngữ ít tài nguyên như tiếng Việt và đủ ba trụ E/S/G. \emph{Hai}, nhiều phương pháp phụ thuộc xếp hạng ESG hoặc nguồn tin bên thứ ba vốn không sẵn có ở thị trường mới nổi. \emph{Ba}, các hướng dựa vào bằng chứng nội văn bản chịu rủi ro \emph{đo lường vòng tròn}. Nghiên cứu này lấp các khoảng trống đó bằng một khung đầy đủ E/S/G cho tiếng Việt, không cần dữ liệu ngoài, và thay grounding vòng tròn bằng đo độ cụ thể của cam kết.

%
%
\section{Phương pháp}
\subsection{Khung và định nghĩa}
Bài toán là đo \emph{decoupling giữa talk và substance} trong báo cáo thường niên của ngân hàng thương mại Việt Nam, đủ ba trụ E/S/G. Vì không có walk-data bên thứ ba cho Việt Nam, chúng tôi đo \emph{substance ngay trong văn bản} qua độ cụ thể của cam kết, thay vì đo bằng grounding nội văn bản (vốn vòng tròn). Cụ thể:
\begin{itemize}
    \item \textbf{Talk} = các câu \emph{cam kết} ESG, xác định bằng bộ phân loại Commitment và lọc qua cổng Topic.
    \item \textbf{Substance (proxy nội văn bản)} = cam kết đó \emph{cụ thể} đến đâu: mơ hồ / có hành động được nêu tên / định lượng quy về chính ngân hàng.
\end{itemize}
Pipeline gồm bốn bước: (1) tách báo cáo thành đơn vị ngữ nghĩa $\le 256$ token; (2) gán trụ E/S/G; (3) phân loại cam kết, lọc qua cổng topic để có mẫu số sạch; (4) chấm độ cụ thể ba mức rồi tổng hợp chỉ số.

\subsection{Corpus và đơn vị phân tích}
Corpus gồm các báo cáo thường niên \emph{bằng tiếng Việt} do chính các ngân hàng thương mại Việt Nam công bố công khai trong giai đoạn 2020--2024 (mỗi ngân hàng một báo cáo/năm).\footnote{Tệp PDF gốc được tải từ chuyên mục Quan hệ cổ đông trên website chính thức của từng ngân hàng; toàn bộ corpus đã số hoá, dữ liệu translate-train và mã nguồn được công khai để tái lập.} Văn bản được trích bằng OCR rồi làm sạch (loại đầu/chân trang, bảng số rời rạc, ký tự nhiễu). Toàn bộ các bước phân loại và chấm điểm phía sau đều thao tác \emph{trực tiếp trên tiếng Việt}, không qua dịch máy sang tiếng Anh. Thay vì cắt cứng theo cửa sổ token (dễ cắt ngang câu, làm mất thực thể/số và bóp méo độ cụ thể), chúng tôi dùng \emph{semantic-text-splitter}: gói các câu liên tiếp đến ngưỡng $\le 256$ token (đo bằng tokenizer Qwen3) theo ranh giới câu. Toàn bộ corpus thu được 13\,812 đơn vị, độ dài trung vị 228 token (p90 $=252$, max $=256$), không có đơn vị rỗng. Mỗi đơn vị mang tối đa một vai trò ngữ nghĩa, giúp tín hiệu phân loại không mâu thuẫn trong cùng đơn vị.

\subsection{Phân loại trụ ESG}
Mỗi đơn vị (bằng tiếng Việt) được gán nhãn đa nhãn vào ba trụ Môi trường (E), Xã hội (S), Quản trị (G); đơn vị không thuộc trụ nào bị loại khỏi phân tích. Chúng tôi tinh chỉnh PhoBERT \cite{nguyen2020phobert} --- mô hình ngôn ngữ được tiền huấn luyện chuyên cho \emph{tiếng Việt} trên kho văn bản 20GB --- với ba đầu ra sigmoid độc lập; nhờ vậy việc phân loại được thực hiện hoàn toàn trên văn bản tiếng Việt gốc, tránh sai lệch do dịch máy. Nhãn huấn luyện thu được theo chiến lược \emph{translate-train}: dịch các tập nhãn chủ đề ESG tiếng Anh (ESGBERT) sang tiếng Việt rồi tinh chỉnh PhoBERT trên miền đích tiếng Việt. Đầu ra của bước này là cờ \texttt{is\_env}, \texttt{is\_soc}, \texttt{is\_gov} cho mỗi đơn vị, đóng vai trò \emph{cổng topic} cho mẫu số ở bước sau.

\subsection{Phân loại cam kết và cổng topic}
Một bộ phân loại nhị phân (CommitmentHF, nền PhoBERT) xác định đơn vị có phải \emph{cam kết} ESG hay không. Bộ phân loại cam kết có xu hướng bắn rộng (nhiều câu ``cam kết'' không gắn trụ ESG nào, ví dụ kết quả tài chính). Để có mẫu số sạch, một đơn vị chỉ vào mẫu số khi \emph{vừa} là cam kết \emph{vừa} có ít nhất một trụ ESG dương. Tỉ lệ đơn vị non-ESG bị loại được ghi nhận minh bạch trong nhật ký chẩn đoán (\texttt{info\_check.json}).

\subsection{Chấm độ cụ thể ba mức}
Trên mỗi đơn vị cam kết, một mô hình ngôn ngữ lớn (Qwen3-1.7B) đóng vai \emph{bộ chấm theo rubric}: phân rã cam kết thành các \emph{item} hành động, rồi suy ra mức cụ thể theo Bảng~\ref{tab:rubric}. Rubric \emph{nhận biết quy-về-chủ-thể}: chỉ tính Mức~2 khi con số là đại lượng đo được \emph{quy về chính ngân hàng}; các con số của NHNN/Chính phủ/toàn ngành/quốc gia hay tiêu chuẩn (ISO, Basel, VIETGAP) \emph{không} được tính là định lượng của chủ thể. Một bước \texttt{verify\_rubric} kiểm tra đầu ra để chống hiện tượng mô hình ``bịa số'' không có trong văn bản.

\begin{table}[!t]
\centering
\caption{Thang độ cụ thể ba mức của cam kết ESG (ví dụ mang tính minh hoạ).}
\label{tab:rubric}
\begin{tabular}{cp{2.6cm}p{4.2cm}p{3.4cm}}
\toprule
Mức & Tên & Định nghĩa & Ví dụ minh hoạ \\
\midrule
0 & Mơ hồ & chỉ khẩu hiệu/tính từ, không kiểm chứng được & ``hướng tới một tương lai xanh và bền vững'' \\
1 & Hành động có tên & có hành động/công cụ/chương trình có tên, quy về chủ thể, nhưng chưa định lượng & ``triển khai hệ thống quản lý carbon nội bộ'' \\
2 & Định lượng & có đại lượng đo được, quy về chính ngân hàng & ``dư nợ tín dụng xanh đạt 74.177 tỷ đồng'' \\
\bottomrule
\end{tabular}
\end{table}

\subsection{Hệ chỉ số CTI / NAR / QDR}
Với mỗi ô (ngân hàng $b$, năm $y$), gọi $N$ là số đơn vị cam kết có ít nhất một trụ ESG dương. Ba chỉ số được định nghĩa từ phân bố mức cụ thể:
\begin{equation}
\mathrm{CTI}(b,y) = \frac{\#\{\text{Mức }0\}}{N}, \quad
\mathrm{NAR}(b,y) = \frac{\#\{\text{Mức }1\}}{N}, \quad
\mathrm{QDR}(b,y) = \frac{\#\{\text{Mức }2\}}{N},
\end{equation}
với $\mathrm{CTI}+\mathrm{NAR}+\mathrm{QDR}=1$. CTI là trục \emph{nói suông} (washing); QDR là trục \emph{thực chất}; NAR là vùng xám ``nói việc cụ thể nhưng chưa có số''. Khác với chỉ số nhị phân, thang ba mức tách được dạng washing tinh vi khi CTI thấp nhưng QDR cũng thấp. Mỗi chỉ số kèm khoảng tin cậy 95\% bằng \emph{percentile bootstrap} (lấy lại mẫu $B=2000$ lần trên các đơn vị trong ô).

\subsection{Công bố chọn lọc}
Để đo cherry-picking, chúng tôi tính tỉ trọng câu ESG theo trụ trên mỗi (ngân hàng, năm):
\begin{equation}
\mathrm{share}(p) = \frac{n_p}{n_E + n_S + n_G}, \quad p \in \{E, S, G\}.
\end{equation}
Lệch mạnh khỏi mặt bằng ngành cùng năm là dấu hiệu tập trung trụ dễ và né trụ khó. Đây là chỉ số \emph{mô tả}, không gộp vào CTI.

%
%
\section{Kết quả và thảo luận}
Khung đo được áp dụng cho 6 ngân hàng (Agribank, BIDV, SHB, Techcombank, Vietcombank, VietinBank) trong 5 năm 2020--2024, tạo thành 30 quan sát panel.\footnote{Đây là tập con đã chạy xong của panel mục tiêu gồm 10 ngân hàng; các ngân hàng còn lại đang được xử lý và sẽ được bổ sung.}

\subsection{Đánh giá bộ phân loại}
Hiệu năng của bộ phân loại trụ ESG và bộ phân loại cam kết được đánh giá trên một Gold Set gồm các câu ESG do người gán nhãn, trích từ báo cáo ngân hàng Việt Nam; cận trên hiệu năng được ước lượng trên gold tiếng Anh theo thiết kế translate-train. Riêng bước chấm độ cụ thể không phải bộ phân loại huấn luyện mà là rubric của mô hình ngôn ngữ lớn, nên được đánh giá bằng \emph{audit thủ công}: lấy mẫu cân bằng theo trụ, người gán tay mức 0/1/2 rồi đo độ đồng thuận với pipeline.
% TODO: điền số liệu Precision/Recall/F1 từ Gold Set khi hoàn tất đánh giá; bảng dưới để khung.
\begin{table}[!t]
\centering
\caption{Hiệu năng bộ phân loại trên Gold Set (đang hoàn thiện đánh giá).}
\label{tab:classifier}
\begin{tabular}{llccc}
\toprule
Tác vụ & Lớp & Precision & Recall & F1 \\
\midrule
\multirow{3}{*}{Trụ ESG}
& E (Môi trường) & -- & -- & -- \\
& S (Xã hội)     & -- & -- & -- \\
& G (Quản trị)   & -- & -- & -- \\
\midrule
Cam kết & Commitment (nhị phân) & -- & -- & -- \\
\midrule
Độ cụ thể & Đồng thuận audit (accuracy/$\kappa$) & \multicolumn{3}{c}{--} \\
\bottomrule
\end{tabular}
\end{table}

\subsection{RQ1 --- Mức độ phổ biến của nói suông}
Trên toàn panel, trung bình $\mathrm{CTI}=0{,}366$ (CI~95\%~$[0{,}334;\,0{,}397]$), $\mathrm{NAR}=0{,}304$ và $\mathrm{QDR}=0{,}330$: hơn một phần ba cam kết ESG là nói suông, chỉ khoảng một phần ba đạt mức định lượng. Bảng~\ref{tab:by_bank} cho thấy khác biệt rõ giữa các ngân hàng: CTI dao động từ $0{,}289$ (Vietcombank, Agribank) tới $0{,}431$ (VietinBank). Vietcombank có QDR cao nhất ($0{,}445$) nhưng số cam kết trung bình mỗi báo cáo thấp nhất ($\approx 30$), nên các ô năm đầu kỳ có khoảng tin cậy rất rộng và cần diễn giải thận trọng.

\begin{table}[!t]
\centering
\caption{CTI/NAR/QDR trung bình theo ngân hàng (2020--2024), khoảng tin cậy bootstrap 95\% của CTI, và số cam kết ESG trung bình mỗi báo cáo ($\bar{N}$).}
\label{tab:by_bank}
\begin{tabular}{lccccc}
\toprule
Ngân hàng & CTI & CTI (CI 95\%) & NAR & QDR & $\bar{N}$ \\
\midrule
Agribank    & 0,293 & [0,237; 0,353] & 0,349 & 0,358 & 48,2 \\
BIDV        & 0,373 & [0,337; 0,398] & 0,326 & 0,301 & 89,0 \\
SHB         & 0,416 & [0,353; 0,460] & 0,283 & 0,302 & 86,0 \\
Techcombank & 0,397 & [0,338; 0,472] & 0,275 & 0,328 & 76,6 \\
Vietcombank & 0,289 & [0,211; 0,373] & 0,266 & 0,445 & 30,2 \\
VietinBank  & 0,431 & [0,399; 0,462] & 0,326 & 0,244 & 82,4 \\
\midrule
Toàn panel  & 0,366 & [0,334; 0,397] & 0,304 & 0,330 & 68,7 \\
\bottomrule
\end{tabular}
\end{table}

\subsection{RQ2 --- Công bố chọn lọc: trụ Môi trường bị né}
Tỉ trọng công bố theo trụ lệch rõ và nhất quán: trung bình $\mathrm{share}(E)=0{,}239$, thấp hơn đáng kể so với $\mathrm{share}(S)=0{,}363$ và $\mathrm{share}(G)=0{,}397$. Kiểm định Friedman trên 30 khối (ngân hàng, năm) bác bỏ mạnh giả thuyết ba trụ được công bố ngang nhau ($\chi^2 = 29{,}9$, $p = 3{,}3\times 10^{-7}$). Hậu kiểm Wilcoxon ghép cặp xác nhận trụ Môi trường thấp hơn cả trụ Quản trị ($p = 8{,}0\times 10^{-6}$) lẫn trụ Xã hội ($p = 1{,}0\times 10^{-6}$), trong khi khác biệt giữa Xã hội và Quản trị không có ý nghĩa ($p = 0{,}064$). Mọi ngân hàng trong mẫu đều công bố trụ Môi trường ít nhất (Bảng~\ref{tab:disclosure}). Đây là bằng chứng định lượng cho hành vi \emph{né chủ đề khó}: ngân hàng dồn ngôn ngữ vào Quản trị và Xã hội --- nơi dễ kể câu chuyện tích cực --- và giảm nhẹ trụ Môi trường, vốn đòi hỏi số liệu phát thải/tín dụng xanh khó kiểm chứng.

\begin{table}[!t]
\centering
\caption{Tỉ trọng công bố trung bình theo trụ trên mỗi ngân hàng. Trụ Môi trường thấp nhất ở toàn bộ ngân hàng.}
\label{tab:disclosure}
\begin{tabular}{lccc}
\toprule
Ngân hàng & share(E) & share(S) & share(G) \\
\midrule
Agribank    & 0,286 & 0,500 & 0,213 \\
BIDV        & 0,267 & 0,321 & 0,411 \\
SHB         & 0,256 & 0,363 & 0,381 \\
Techcombank & 0,179 & 0,385 & 0,436 \\
Vietcombank & 0,243 & 0,298 & 0,459 \\
VietinBank  & 0,202 & 0,314 & 0,484 \\
\midrule
Trung bình  & 0,239 & 0,363 & 0,397 \\
\bottomrule
\end{tabular}
\end{table}

\subsection{RQ3 --- Khoảng cách thực chất theo thời gian}
Phân tích theo năm (Bảng~\ref{tab:by_year}) cho thấy một nghịch lý. Khối lượng cam kết ESG tăng mạnh và có ý nghĩa thống kê theo thời gian (Spearman $\rho = 0{,}59$ giữa $\bar{N}$ và năm, $p < 0{,}001$): trung bình mỗi báo cáo tăng từ $\approx 45$ cam kết năm 2020 lên $\approx 101$ năm 2024. Tuy nhiên, \emph{cơ cấu} độ cụ thể gần như đứng yên: cả CTI ($\rho = 0{,}02$, $p = 0{,}90$) lẫn QDR ($\rho = -0{,}14$, $p = 0{,}47$) đều không có xu hướng thời gian có ý nghĩa. Nói cách khác, các ngân hàng \emph{nói nhiều hơn về ESG} nhưng \emph{không cụ thể hơn}: phần tăng thêm chủ yếu là cam kết, chứ không phải nội dung định lượng.

Tương quan thành phần củng cố diễn giải này. Số cam kết càng lớn thì QDR càng thấp (Spearman $\rho = -0{,}55$, $p = 0{,}002$) và CTI càng cao ($\rho = 0{,}36$, $p = 0{,}048$): việc mở rộng khối lượng công bố đi kèm \emph{pha loãng} nội dung thực chất. Phát hiện này cũng đính chính một nhận định trong dòng nghiên cứu trước đó về một ``đỉnh rủi ro năm 2022 rồi cải thiện'': với thước đo độ cụ thể, dữ liệu không ủng hộ một quỹ đạo cải thiện đơn điệu --- CTI thậm chí cao nhất ở năm 2023.

\begin{table}[!t]
\centering
\caption{CTI/NAR/QDR trung bình theo năm (khoảng tin cậy bootstrap 95\% của CTI lấy qua 6 ngân hàng) và số cam kết ESG trung bình mỗi báo cáo ($\bar{N}$).}
\label{tab:by_year}
\begin{tabular}{lccccc}
\toprule
Năm & CTI & CTI (CI 95\%) & NAR & QDR & $\bar{N}$ \\
\midrule
2020 & 0,377 & [0,279; 0,475] & 0,283 & 0,340 & 44,8 \\
2021 & 0,317 & [0,236; 0,395] & 0,301 & 0,383 & 47,5 \\
2022 & 0,377 & [0,338; 0,413] & 0,323 & 0,299 & 65,0 \\
2023 & 0,404 & [0,352; 0,444] & 0,284 & 0,312 & 85,0 \\
2024 & 0,356 & [0,318; 0,395] & 0,330 & 0,314 & 101,3 \\
\bottomrule
\end{tabular}
\end{table}

\subsection{Kiểm định độ ổn định xếp hạng}
Để kiểm tra liệu CTI có dùng để \emph{xếp hạng} ngân hàng một cách đáng tin hay không, chúng tôi lấy lại mẫu bootstrap các cam kết trong từng (ngân hàng, năm), tính lại CTI cấp ngân hàng và đo Kendall $\tau$ của thứ hạng so với gốc. Thứ hạng ổn định ở mức khá: $\tau$ trung bình $=0{,}67$ (phân vị 5\% $=0{,}20$), và ngân hàng rủi ro cao nhất (VietinBank) được giữ ở đỉnh trong 78\% lần lấy mẫu. Kiểm định bỏ-một-năm cho $\tau$ tối thiểu $=0{,}73$, cho thấy thứ hạng không bị chi phối bởi bất kỳ năm đơn lẻ nào. Độ ổn định ở mức vừa phải (chưa cao tuyệt đối) phản ánh đúng thực tế cỡ mẫu mỗi ngân hàng còn nhỏ và các giá trị CTI khá sát nhau --- một giới hạn sẽ giảm khi mở rộng panel.

\subsection{Thảo luận}
Ba phát hiện bổ sung cho nhau thành một bức tranh nhất quán về ESG-washing trong ngành ngân hàng Việt Nam: nói suông phổ biến (RQ1), trụ Môi trường bị né một cách hệ thống (RQ2), và sự bành trướng cam kết không đi kèm cải thiện thực chất (RQ3). Quan trọng hơn, các kết luận này được rút ra \emph{mà không cần} dữ liệu bên thứ ba và \emph{không} dựa vào cấu trúc grounding vòng tròn --- cho thấy đo độ cụ thể nội văn bản là một thước đo khả thi và trung thực cho các thị trường thiếu dữ liệu.

\paragraph{Vì sao tín hiệu nội văn bản vẫn có giá trị.} Một phản biện chính đáng là phân tích chỉ dùng bản thân báo cáo nên không nắm được yếu tố bên ngoài hay khoảng cách giữa lời viết và hành động thực tế. Chúng tôi xác định rõ đây là \emph{phạm vi có chủ đích} chứ không phải thiếu sót bị bỏ qua, và lập luận tín hiệu nội văn bản vẫn mang thông tin nhờ lý thuyết \emph{cheap talk} \cite{farrell1996cheap} và tín hiệu tốn chi phí. Một cam kết mơ hồ (``hướng tới tương lai xanh'') gần như không tốn chi phí và không thể bị bác bỏ; ngược lại, một cam kết \emph{định lượng, quy về chính ngân hàng} (``dư nợ tín dụng xanh đạt X tỷ'') là phát ngôn \emph{có thể chứng minh sai}, tạo ràng buộc pháp lý và rủi ro uy tín. Do đó, \emph{phân bố độ cụ thể} của cam kết phản ánh mức độ tổ chức sẵn sàng chịu trách nhiệm --- một chiều của decoupling đo được ngay trong văn bản, độc lập với việc có dữ liệu kiểm chứng bên ngoài hay không. Khung của chúng tôi vì thế không thay thế kiểm toán ESG bên thứ ba mà \emph{bổ sung} nó: ở những thị trường chưa có walk-data, đo độ cụ thể là tầng sàng lọc đầu tiên, có thể mở rộng quy mô, để khoanh vùng báo cáo và trụ ESG đáng nghi trước khi tốn nguồn lực điều tra bên ngoài.

%
%
\section{Kết luận}
Nghiên cứu đề xuất một khung NLP đo rủi ro ESG-washing trong báo cáo thường niên của ngân hàng thương mại Việt Nam qua \emph{độ cụ thể của cam kết}, thay cho các cấu trúc phụ thuộc dữ liệu bên thứ ba hay bằng chứng nội văn bản mang tính vòng tròn. Bằng cách phân loại trụ ESG, nhận diện cam kết và chấm độ cụ thể ba mức, khung tổng hợp ba chỉ số bổ sung nhau CTI/NAR/QDR. Áp dụng cho 6 ngân hàng giai đoạn 2020--2024, kết quả cho thấy nói suông chiếm trung bình 36{,}6\% cam kết, trụ Môi trường bị công bố thiếu có ý nghĩa thống kê, và khối lượng cam kết tăng nhanh nhưng tỉ trọng thực chất không cải thiện; thứ hạng rủi ro giữa các ngân hàng ổn định dưới kiểm định bootstrap.

Nghiên cứu có một số hạn chế cần ghi nhận minh bạch. \emph{Một}, độ cụ thể là proxy nội văn bản chứ không phải bằng chứng thực thi bên ngoài; không có walk-data nên không thể đo trực tiếp recall/precision của washing. \emph{Hai}, nhãn cam kết/độ cụ thể bắt nguồn từ dữ liệu khí hậu tiếng Anh được translate-train, có thể chưa chuyển giao hoàn hảo sang trụ S/G tiếng Việt. \emph{Ba}, kết quả hiện dựa trên 6 ngân hàng. Hướng phát triển gồm: mở rộng panel lên đủ 10 ngân hàng, hoàn tất audit thủ công để báo cáo độ đồng thuận của thang độ cụ thể, và bổ sung các kiểm định nhiễu (chèn cam kết mơ hồ, kiểm tra chống lối tắt ``có chữ số là cụ thể'') nhằm củng cố tính bền vững của thước đo.

\begin{credits}
\subsubsection{\discintname}
Các tác giả tuyên bố không có xung đột lợi ích liên quan đến nội dung bài báo.
\end{credits}

%
% ---- Tài liệu tham khảo ----
\bibliographystyle{splncs04}
\bibliography{references}

\end{document}
